\subsection{All marked symmetric quotients at once}

\begin{statement}{D.41}{Proved}{A simultaneous symmetric extension and an explicit remaining projective-linear kernel}
Let
\[
 \mathscr S=\{M\lhd B:[B:M]<\infty,\ B/M\simeq S_n
                         \text{ for some }n\ge5\}.
                                                               \tag{D.41.1}
\]
There is one coordinate for each actual marked kernel $M$.
Put $S_M=B/M$, $A_M=S_M'$, and let $\epsilon_M:S_M\to C_2$
be its sign map. The actual simultaneous source image is
\[
 \mathcal S
   =\left\{(g_M)\in\prod_{M\in\mathscr S}S_M:
             \epsilon_M(g_M)\text{ is independent of }M\right\}.
                                                               \tag{D.41.2}
\]
The map $\widehat B\to\mathcal S$ is onto, and
\[
 \ker(\mathcal S\to C_2)=\prod_{M\in\mathscr S}A_M,\qquad
 \mathcal S\simeq
       \left(\prod_{M\in\mathscr S}A_M\right)\rtimes C_2.
                                                               \tag{D.41.3}
\]
The section in this semidirect product is the actual
tuple of the images of $s$, and its action on each
$A_M$ is the indicated conjugation. Every nonabelian
finite quotient of $\mathcal S$ has some $S_n$ quotient.

Retain $\mathcal U=\mathcal T\times_{C_6}\mathcal R$ and
the exact kernels $R_T,K_{\mathcal R}$ of D.40. With
$K_{\mathcal S}=\bigcap_{M\in\mathscr S}\overline M$,
there is an exact marked isomorphism
\[
 \mathcal W
   :=\widehat B/(R_T\cap K_{\mathcal R}\cap K_{\mathcal S})
     \simeq\mathcal U\mathbin{\times_{C_2}}\mathcal S.
                                                               \tag{D.41.4}
\]
The automorphism
$(\alpha_{39},1_{\mathcal R},1_{\mathcal S})$ is an
actual continuous peripheral automorphism of
$\mathcal W$, with multiplier one and both hexagons.
Its simultaneous conjugators and global gauge are
specified below. It extends D.40 exactly and acts
as the identity on every actual full symmetric
quotient in (D.41.1), including both displayed
$S_7$ and $S_{10}$ quotients.

Let $M_7$ be the kernel of D.40's specified $S_7$ map
and let $\mathcal V$ be its source quotient there.
The canonical projections have kernels
\[
 \begin{split}
 \ker(\mathcal W\to\mathcal V)
     &=\prod_{M\ne M_7}A_M,\\
 \ker(\mathcal W\to\mathcal T)
     &=\mathcal R'\times\prod_M A_M .
 \end{split}                                             \tag{D.41.5}
\]
The products carry their profinite topology, and
$\mathcal R'$ retains its closed-commutator definition
from D.40.

An explicit remaining finite source map is
\[
 \begin{split}
 \sigma:\widehat B&\twoheadrightarrow
                   J=\operatorname{PGL}_2(\mathbb F_7),\\
 s&\longmapsto
      \left[\begin{pmatrix}1&0\\0&-1\end{pmatrix}\right],
 \qquad
 r\longmapsto
      \left[\begin{pmatrix}1&-1\\1&0\end{pmatrix}\right].
 \end{split}                                             \tag{D.41.6}
\]
The remaining full-source kernel has exact image
\[
 \sigma(R_T\cap K_{\mathcal R}\cap K_{\mathcal S})
       =\operatorname{PSL}_2(\mathbb F_7).                \tag{D.41.7}
\]
Thus it is nontrivial. The canonical good-prime
congruence map $\rho_7:B\to\operatorname{PSL}_2(\mathbb F_7)$
and (D.41.6) have actual joint images
\[
 \operatorname{im}(\rho_7,\sigma)
      =\operatorname{PSL}_2(\mathbb F_7)
             \times\operatorname{PGL}_2(\mathbb F_7),
 \quad
 \operatorname{im}(\rho_7,\sigma)|_{B_{\rm even}}
      =\operatorname{PSL}_2(\mathbb F_7)^2.               \tag{D.41.8}
\]
Here $B_{\rm even}$ is the kernel of the original sign
map $s\mapsto1,r\mapsto0$. In particular the kernel in
(D.41.7) kills the canonical $\rho_7$ coordinate while
surjecting under the specified second marked map.

The source $\widehat F$ maps onto
$\ker(\mathcal W\to S_3)$ with the same displayed
kernel. D.39's compact-cover matrices and fixed-fibre
CM pair are unchanged. D.45 excludes a full profinite lift
of this recorded action. Formula (D.41.7) describes the
source kernel, while the whole prescribed CM pair remains
open for a replacement action.

This addresses Part D(e), section 51.16, folio 734 verso,
scan 174, and D-after-LM2(1), section 51.15, folios
725--727, scans 156--160.
\end{statement}
\begin{proof}
\emph{The exact simultaneous symmetric image.}
For every epimorphism $B\to S_n$, $n\ge5$, the image
of $r$ is even because its order divides three.
The image of $s$ must be odd, since otherwise the
whole image would lie in $A_n$. Thus all the
characters $\epsilon_M$ pull back to the same
original sign character. In particular
$M\subset B_{\rm even}$ for every $M\in\mathscr S$.
The restriction $B_{\rm even}\to A_M$ is onto
and has kernel exactly $M$. Restrictions coming
from two different members of $\mathscr S$
therefore have different kernels, even when
their abstract alternating groups have the
same degree.

These restrictions are simultaneously onto every
finite product of their simple targets. To prove
this, adjoin one target at a time. A subdirect
product of a previous full product and a new
nonabelian simple group either is the full
product or identifies a common simple quotient.
Every epimorphism from a product of nonabelian
simple groups to a simple group factors through
one coordinate: the images of the factors are
commuting normal subgroups, and two nontrivial
images would make the target abelian.
A common simple quotient would therefore
identify the new restriction kernel with one
of the preceding kernels, which was excluded.
This proves finite-product surjectivity.
Compactness gives
\[
 \overline{B_{\rm even}}\twoheadrightarrow
                  \prod_{M\in\mathscr S}A_M.            \tag{D.41.9}
\]

Multiplication by the actual image of $s$ supplies
every tuple with common odd sign. This proves
(D.41.2) and (D.41.3), including inverse maps:
write a tuple with common sign $\epsilon$ as
\[
 (g_M)=(a_Ms_M^\epsilon),\qquad
 a_M=g_Ms_M^{-\epsilon}\in A_M,
 \quad s_M=sM.                                         \tag{D.41.10}
\]
The action in the semidirect product is therefore
literally conjugation by each $s_M$.

This construction also explains repeated
epimorphisms with the same kernel. Their quotient
maps differ by the unique induced automorphism
of the quotient, so their image is the graph of
that automorphism. Dropping repeated coordinates
and reinserting them by these maps are mutually
inverse marked isomorphisms. In a finite family
with repetitions, the even image is correspondingly
a product of diagonal simple blocks. The normalizer
of each such block inside the product of its
alternating factors is the block itself: a
normalizing tuple induces the same conjugation
in every identified factor, so its differences
centralize the centreless simple group and are
trivial. This proves the normalizer assertion
directly, and also gives the equality of the
even subgroup and the commutator subgroup of
each finite simultaneous image.

\emph{Finite quotients and the common sign.}
Consider a finite coordinate image
$H=D\rtimes\langle s\rangle$, where
$D=\prod_{i=1}^j A_{M_i}$. A normal subgroup of
$D$ is a product of some of its simple factors.
Indeed, a nontrivial projection to a factor,
commuted with that factor, gives a nontrivial
normal subgroup of that factor and hence
contains it; the trivial centre removes every
remaining projection.

Let $N\lhd H$. If $N$ contains an odd element
$a s$, then on every factor its conjugation
is given by an odd permutation. This action
cannot be the identity: the centralizer of
$A_n$ in $S_n$ is trivial, as proved in D.40.
Consequently $[N,A_{M_i}]$ is nontrivial and
contains $A_{M_i}$ for every $i$. Thus $N$
contains $D$ and an odd element, so $N=H$.
Every proper normal subgroup of $H$ is therefore
a product of some alternating factors.
A nonabelian quotient retains at least one
factor and maps onto its corresponding
$S_{n_i}$. The only other nontrivial quotient
is the common $C_2$.

Every finite quotient of $\mathcal S$ factors
through a finite coordinate image. The preceding
argument proves the assertion about all its
nonabelian finite quotients, and also
$\mathcal S^{\rm ab}=C_2$.

The argument of D.40 excluding an $S_7$ quotient
of $\mathcal U$ applies to every $S_n$, $n\ge5$.
In a finite source image
$J_0\subset T_0\times G_1\times\cdots\times G_k$,
$T_0$ is solvable and the $G_i$ are perfect.
Under a proposed epimorphism to $S_n$, the
$T_0$-coordinate kernel has image containing
$A_n$. Each perfect-coordinate kernel has
image $S_n$: the derived subgroup maps onto
that coordinate, so the kernel maps onto the
sign quotient and is normal. Their iterated
commutator is trivial in the coordinate
product but has image containing the
nontrivial perfect group $A_n$.

It follows that every common finite quotient
of $\mathcal U$ and $\mathcal S$ is a quotient
of the common sign $C_2$. A nonabelian common
profinite quotient would have a nonabelian
finite quotient and hence an $S_n$ quotient,
which has just been excluded. The maximal
common marked quotient is exactly $C_2$.
The closed subdirect-product construction of
D.40 now proves (D.41.4) in both directions,
including its intersection of source kernels.

\emph{The actual compatible automorphism.}
The normalized D.39 conjugators have zero
$B$-abelianization. Pair them with the identity
in every $\mathcal R$ and $\mathcal S$ coordinate.
They agree on the common $C_6$ and $C_2$ quotients,
so are actual elements of the fibre product.
Before normalization the coordinate conjugators
are $x^{-k}$, using the actual single
$k\in\widehat{\mathbb Z}$ of D.39. Their coordinate
actions are $\operatorname{Int}(x^{-k})$, which
fixes the cusp. The common gauge by $x^k$
changes every additional conjugator to one,
and every additional action to the identity.
The resulting commutator is
$(f_{39},1_{\mathcal R},1_{\mathcal S})$.

The source kernel in (D.41.4) is contained in
$R_T\subset\widehat F'$, so the marked
abelianization is still $\widehat{\mathbb Z}^2$.
The same gauge calculation gives a lift of
this commutator in $\widehat F'$.
The component automorphisms and inverses
preserve their common characters, and hence
give an actual automorphism of the fibre
product. The torsion and cusp equations
give both hexagons, compatibly at every
finite level. No new multiplier is introduced.

Projection to the $M_7$ coordinate of
$\mathcal S$ recovers D.40's action exactly.
The kernel of that projection has even sign
and consists of the independent remaining
alternating factors. Projection to
$\mathcal T$ additionally sets the $C_6$
image of $\mathcal R$ to zero. These
calculations prove both kernels in
(D.41.5). In particular the actual $S_{10}$
coordinate now has the identity action;
D.40's $A_{10}$ layer is included in the
new source quotient.

\emph{An explicit remaining marked quotient.}
The program works in projective matrices
over $\mathbb F_7$, choosing the representative
whose first nonzero entry is one.
The two matrices in (D.41.6) satisfy the
projective relations $s^2=r^3=1$ and generate
all $336$ elements. Their $t=sr$ has order eight
and $x=t^{-2}$ has order four. The square determinant
subgroup
$L=\operatorname{PSL}_2(\mathbb F_7)$ has
order $168$, and the determinant square-class
map gives $J/L=C_2$. The group $B_{\rm even}$
is generated by $r,srs$; their images generate
all $168$ elements of $L$.

For a complete elementary simplicity check,
the program exhausts the conjugacy classes
of $L$. Their sizes are
\[
 1,\quad21,\quad42,\quad56,\quad24,\quad24.
                                                               \tag{D.41.11}
\]
The subgroup generated by every nonidentity
class has order $168$. Every nontrivial
normal subgroup therefore equals $L$.
The group is nonabelian, so is perfect.
The centralizer of $L$ in $J$ has order
one, checked on the generating pair
$r,srs$. It follows, by intersecting a
normal subgroup with $L$, that the normal
subgroups of $J$ are exactly $1,L,J$.
In particular $J'=L$.

The action on the eight points of
$\mathbb P^1(\mathbb F_7)$ is faithful, and
its permutation sign equals determinant
square class. Its image has order $336$,
whereas $S_8$ has order $40320$. Thus this
specific permutation representation is not
an epimorphism to $S_8$ and is not one of
the full symmetric coordinates in (D.41.1).
This follows from its explicit action map
and image order.

There is no quotient $\mathcal W\to J$.
A proposed continuous epimorphism factors
through a finite marked image
\[
 H\subset T_0\times\prod_{i=1}^a G_i
                       \times\prod_{j=1}^b S_{n_j},
                                                               \tag{D.41.12}
\]
with $T_0$ solvable and each $G_i$ perfect.
The $T_0$-coordinate kernel has image containing
$L$, since its quotient is solvable.
Each perfect-coordinate kernel has image
$J$: the equality $H=H'N_i$ forces its normal
image to meet the nontrivial determinant
coset. Each symmetric-coordinate kernel
also has image containing $L$. Otherwise
its image would be trivial, and $J$ would
be a quotient of $S_{n_j}$. The only
nontrivial proper quotient of that group
is $C_2$, and $336$ is not a factorial
$n!$ for $n\ge5$, so that is impossible.
The iterated commutator of all coordinate
kernels is trivial in (D.41.12), while its
image contains the nontrivial perfect
group $L$. This contradiction proves the claim.

The image of the kernel in (D.41.7) is normal
in $J$ and lies in $J'=L$, since the kernel
is contained in $\widehat F'\subset\widehat B'$.
It cannot be trivial, as that would make
$\sigma$ factor through $\mathcal W$.
Simplicity of $L$ now proves (D.41.7).

Finally the canonical $\rho_7$ has the
different, previously fixed modular matrices
\[
 \rho_7(s)=
 \left[\begin{pmatrix}0&1\\-1&0\end{pmatrix}\right],
 \qquad
 \rho_7(r)=
 \left[\begin{pmatrix}0&-1\\1&-1\end{pmatrix}\right].
                                                               \tag{D.41.13}
\]
The program evaluates the paired generators
from (D.41.6) and (D.41.13), obtaining all
$168\cdot336=56448$ paired elements.
Using the actual even-source generators
$r,srs$ gives all $168^2=28224$ paired
elements. These full counts prove
(D.41.8) under the displayed maps.
The canonical $\rho_7$ kernel occurs among the
subgroups intersected in $K_{\mathcal R}$,
whereas (D.41.7) remains
onto under the second map. The two marked
occurrences of the simple group are thus
compared explicitly. The source kernel and
its open $\widehat F$ restriction are exactly
those in (D.41.4); full profinite realization
has not been inferred from the quotient action.
\end{proof}
